\scanpage{96}
\mspage{47}
given $G$, the set $\widetilde A(G)$ and $I(G)$ (both equipped with an \inserted{``origin''} map with values in $S(G)$).

For this, we \unsure{have|to|will} use the notion of ``graphs with flattened branches'', which generaliz\unsure{es} that of a graph, and topologically corresponds to the classification of \unsure{\deleted{4}|4|\deleted{q}|q} pairs $(X,S)$ of a compact topological space $X$, and a finite subset $S$, such that a) $X$ is a star \unsure{in a|in a} neighborhood of $\forall$ point of $S$, \deleted{\unsure{and|a}} b) \deleted{\unsure{S}} every connected component of $X\setminus S$ is a noncompact $1$-manifold with boundary (for we want every \deleted{\unsure{closure} of} ``edge'' to have a vertex). \unsure{It|It} amounts to the same to say that $X$ is a topological $1$-complex, and $S$ a finite subset which contains the vertices of order $\geq3$, and such that no connected component of $X\setminus S$ is compact (i.e. that every connected component of $X$ contains a vertex, i.e. $S\unsure{\rightharpoondown|{-}}\unsure{\Pi|\pi}_0(X)$ is surjective).

This is classified combinatorially by systems $(S,\unsure{\widehat A|\widehat I},\sigma_{\unsure{\widehat A|\widehat I}},o_{\widetilde A})$, where \deleted{$\unsure{\widetilde A=\sigma_{\widetilde A}|\widetilde A,\sigma_A|\widetilde A,o_A}$} $S,\unsure{\widehat A|\widehat I}$ are finite sets, $\sigma_{\unsure{\widetilde A|\widehat A}}$ is an involution of $\widetilde A$, and $o_{\widetilde A}:\widetilde A\unsure{\rightharpoondown|{-}}S$ a
\margin{95}
